% !TEX program = xelatex
\documentclass[11pt,a4paper]{article}

% --- 中文支援（XeLaTeX + xeCJK；若用 pdfLaTeX 請改用 CJKutf8） ---
\usepackage{fontspec}
\usepackage[BoldFont, SlantFont]{xeCJK}
% 跨平台中文字型 fallback：依序嘗試 macOS → Windows → Overleaf/Linux，最後一棒一定存在
\IfFontExistsTF{Heiti TC}
  {\setCJKmainfont{Heiti TC}}                       % macOS 內建（黑體-繁）
  {\IfFontExistsTF{Microsoft JhengHei}
    {\setCJKmainfont{Microsoft JhengHei}}           % Windows 內建（微軟正黑體）
    {\setCJKmainfont{Noto Sans CJK TC}}}            % Overleaf / Linux 內建（Google Noto）


\usepackage{amssymb}        % 保險，CVaR 等若用到額外數學符號

% --- 與 TA 範本一致的套件 ---
\usepackage{amsfonts}
\usepackage{amssymb}
\usepackage{amsmath}
\usepackage{amsthm}
\usepackage{epsfig}
\usepackage{graphicx}
\usepackage{natbib}
\usepackage{textcomp}
\usepackage{booktabs}
\usepackage{multirow}
\usepackage{url}
\usepackage{color}
\usepackage{fullpage}
\usepackage[capitalize]{cleveref}
\usepackage{mathtools}
\usepackage{enumitem}
\usepackage{authblk}
\usepackage{algorithm}
\usepackage{algpseudocode}
\usepackage{subcaption}

% --- 版面 ---
\renewcommand{\baselinestretch}{1.25}
\parskip=5pt
\parindent=20pt
\footnotesep=5mm

% --- theorem-like 環境 ---
\newtheorem{lem}{Lemma}
\newtheorem{prop}{Proposition}
\newtheorem{thm}{Theorem}
\newtheorem{defn}{Definition}

\newenvironment{enumerateTight}{\begin{enumerate}\vspace{-8pt}}{\end{enumerate}\vspace{-8pt}}
\newenvironment{itemizeTight}{\begin{itemize}\vspace{-8pt}}{\end{itemize}\vspace{-8pt}}
\leftmargini=25pt
\leftmarginii=12pt

% --- 標題 ---
\title{Operations Research, Spring 2026 (114-2) \\
Final Project Report \\[4pt]
\large 主題型股票投資之兩階段再平衡決策支援系統 \\
\normalsize 結合投資組合最佳化、整數張數轉換與現金流守恆之混合整數規劃模型}

\author{顧懷允 (B13701220)}
\author{林子宸 (B13703048)}
\author{劉威廷 (B13705036)}
\author{闕以諾 (B13303079)}
\affil{Department of Information Management, National Taiwan University}

\date{Group ID：第 12 組}

\begin{document}
\maketitle

\begin{abstract}
\noindent 本研究針對「主題型台股投資的再平衡決策」這個過去由投資經理人徒手處理的流程，建構出一套兩階段混合整數線性規劃（MILP）決策支援系統。第一階段以「投資吸引力分數最大化」為目標，在持股檔數、產業集中度、Beta、CVaR、周轉率與流動性等多重限制下，求解理想連續權重 $w^\star$ 與持股二元變數 $z^\star$；第二階段則把 $w^\star$ 落地為「整數張數、現金守恆、買賣互斥、考量手續費與證交稅」的可下單清單。我們另外提出一個多起點啟發式演算法（greedy + water-filling + repair + 1-swap local search），不呼叫商用 solver 即可在毫秒至秒級給出可行解。以台股 2026/05/20 為再平衡日、1{,}052 檔候選股、240 個歷史情境的真實案例上，啟發式相對 MILP 最佳解之 optimality gap 為 \textbf{1.05\%}、相對 Top-$K$ 等權基準則高出 \textbf{17.71\%}；Stage 2 的整數張數落地之累積偏離為 \textbf{0.0366}、現金剩餘僅 \textbf{0.238\%}。
\end{abstract}



% ============================================================
\section{Introduction}
\label{sec-intro}

\textbf{背景與動機。}{}主題型投資（thematic investing）近年成為台股市場中常見的決策框架——投資人關心的不再只是「個股有沒有漲」，而是「某個主題（如 AI 伺服器、先進封裝、散熱、衛星通訊）會不會延續走強」。然而，\textbf{從一個主題到一份真正可下單的調倉清單}之間仍隔著相當長的決策鏈：先確認哪些個股屬於該主題、再評估每檔股票的報酬潛力與風險、再決定要放多少權重、最後還得處理「台股最小交易單位是一張＝1{,}000 股」「現金不能變負」「買賣手續費與證交稅必須一起算」等實務限制。

\textbf{觀察到的痛點。}{}我們在訪談熟悉台股實務的研究員後歸納出三個常被忽略的決策落差：
\begin{itemizeTight}
    \item \textbf{權重 → 張數的落地問題}。傳統 Markowitz / Black–Litterman 等模型輸出 7.36\%、4.82\% 等連續權重，但實際下單必須是整數張數；越接近最低權重（例如 1\%），離散化造成的偏離越大。
    \item \textbf{現金不平衡}。買賣手續費與證交稅讓「賣 X 元、買 Y 元」事後並非保本；若不在最佳化內處理，現金可能不夠交割。
    \item \textbf{多風險限制的同時滿足}。Beta（市場風險）、CVaR（極端下檔）、產業集中度、單檔流動性、單期周轉率原本由人工反覆嘗試，K=20 檔的問題就足以使人工方案陷入次優或不一致。
\end{itemizeTight}

\textbf{決策者與權衡。}{}本研究的決策者是「持有資金 $V^0$ 的主題型投資組合經理人」。其每一期面對的核心權衡是：
\begin{itemizeTight}
    \item \emph{報酬潛力 vs.\ 風險控制}：高吸引力分數（$\mu$）的股票通常 Beta 較高，硬塞權重會違反 $\beta^{\max}$；極端日報酬大的股票會推升 CVaR。
    \item \emph{理想配置 vs.\ 可執行性}：權重 $w_i^\star$ 越精細，整數張數的離散偏離越難壓低；張數越保守，現金又會大量閒置。
    \item \emph{交易成本 vs.\ 周轉}：頻繁調倉抓信號代價是手續費與滑價；過低的周轉率限制 $\tau$ 又會讓組合難以反應新觀點。
\end{itemizeTight}

\textbf{為什麼這是 OR 問題。}{}上述每個權衡都對應到一條（或一組）線性限制式，決策變數天然含連續權重與整數張數，整個問題剛好落在 MILP 標準型內、可由 Gurobi／CBC 求解。我們因此設計兩個獨立但耦合的模型——Stage 1 解策略層（理想權重）、Stage 2 解執行層（張數＋現金）——並以基準對照驗證其相較簡化策略（Top-$K$ 等權、僅 Stage 1 連續權重）的優勢。



% ============================================================
```latex
\section{Problem description}
\label{sec-problem}

我們把「主題型再平衡」拆成兩個串接的決策問題。第一階段是\textbf{投資組合配置決策}：在每一個再平衡日 $t$，投資組合經理人需從候選股池中選出恰好 $K$ 檔股票，並決定各股票之目標權重 $w_i$。配置結果須滿足權重總和、持股檔數、單股上下限、產業／主題集中度、Beta、流動性、周轉率與 CVaR 風險限制；目標則是最大化 $\mu^\top w$，其中 $\mu_i$ 表示股票 $i$ 的事前吸引力分數。

第二階段是\textbf{交易落地決策}：給定第一階段求得之理想權重 $w^\star$ 與持股名單 $z^\star$，並結合期初張數 $x^0$、期初現金 $C^0$ 與每張價格 $P_i$，決定每檔股票應買入張數 $u_i$、賣出張數 $v_i$、最終持有張數 $x_i$ 與剩餘現金 $C$。此階段要求庫存守恆、現金非負、持股名單延用 Stage 1 結果、買賣互斥，且所有交易張數皆為非負整數。目標是使實際配置盡量貼近 $w^\star$，並加入小幅交易標的數懲罰以避免零碎下單。

\textbf{原本人工 vs.\ 系統化。}
\Cref{tab:auto} 摘要原流程中需人工判斷的環節，以及本研究如何以最佳化模型將其系統化。

\begin{table}[!hbt]
\centering
\scriptsize
\setlength{\tabcolsep}{4pt}
\renewcommand{\arraystretch}{0.9}
\begin{tabular}{p{0.42\textwidth} p{0.50\textwidth}}
\toprule
原本需人工處理之事項 & 本研究之系統化處理方式 \\
\midrule
挑選主題股票與比較投資吸引力 & 依主題、市值、成交值與資料完整性建池，並自動計算可量化 score \\
控制產業或主題集中度 & 將產業權重上限納入最佳化限制式 \\
控制市場風險、下檔風險與換股幅度 & 加入 Beta、CVaR 與周轉率限制 \\
將理想權重轉成可下單張數 & 以整數變數決定買／賣／最終持股張數，並加入現金守恆限制 \\
產生可執行調倉清單 & 輸出 order sheet：買入／賣出張數與最終現金餘額 \\
\bottomrule
\end{tabular}
\caption{人工流程與本研究系統化處理方式}
\label{tab:auto}
\end{table}


% ============================================================
\section{Mathematical model}
\label{sec-model}

我們以 MILP 描述上述兩階段問題。CVaR、周轉率、L1 偏離與買賣互斥皆以線性限制式表示，因此整體可直接交由 MILP 求解器處理。

\subsection{Stage 1：理想投資組合最佳化}

令 $i\in I=\{1,\dots,N\}$ 為候選股，$j\in J=\{1,\dots,M\}$ 為產業／主題群組，$S_j\subseteq I$ 為群組 $j$ 內股票集合，$s\in\Omega$ 為 CVaR 歷史情境。主要參數包含投資吸引力分數 $\mu_i$、Beta $\beta_i$、期初權重 $w_i^0$、單股下／上限 $L,U$、持股檔數 $K$、群組上限 $\gamma_j$、流動性上限 $\ell_i$、周轉率上限 $\tau$、情境報酬 $r_{i,s}$、信賴水準 $\alpha$ 與 CVaR 上限 $\bar C^{\text{risk}}$。決策變數為目標權重 $w_i$、選股變數 $z_i$、周轉率線性化變數 $\delta_i^+,\delta_i^-$、VaR 輔助變數 $\eta$ 與情境超額損失 $\xi_s$。

\textbf{Stage 1 模型：}
{\scriptsize
\setlength{\arraycolsep}{3pt}
\begin{align}
\max_{w,z,\delta^\pm,\eta,\xi}\quad 
& \sum_{i\in I}\mu_i w_i \label{obj1}\\[-1mm]
\text{s.t.}\quad
& \sum_{i\in I}w_i=1, 
&& \sum_{i\in I}z_i=K, \label{c1:basic}\\
& Lz_i\le w_i\le Uz_i,\quad \forall i\in I,
&& w_i\le \ell_i,\quad \forall i\in I, \label{c1:boxliq}\\
& \sum_{i\in S_j}w_i\le \gamma_j,\quad \forall j\in J,
&& \sum_{i\in I}\beta_iw_i\le \beta^{\max}, \label{c1:indbeta}\\
& w_i-w_i^0=\delta_i^+-\delta_i^-,\quad \forall i\in I,
&& \sum_{i\in I}(\delta_i^++\delta_i^-)\le \tau, \label{c1:turn}\\
& \xi_s\ge -\sum_{i\in I}r_{i,s}w_i-\eta,\quad \forall s\in\Omega,
&& \eta+\frac{1}{(1-\alpha)|\Omega|}\sum_{s\in\Omega}\xi_s\le \bar C^{\text{risk}}, \label{c1:cvar}\\
& w_i,\delta_i^\pm,\xi_s\ge0,\quad \eta\in\mathbb R,
&& z_i\in\{0,1\},\quad \forall i\in I. \nonumber
\end{align}
}

其中 \Cref{c1:turn} 將 $|w_i-w_i^0|$ 拆為正負兩個非負變數，以線性化周轉率限制；\Cref{c1:cvar} 為 Rockafellar--Uryasev CVaR 線性化形式。

\subsection{Stage 2：交易落地與整數張數轉換}

給定 Stage 1 的解 $(w^\star,z^\star)$，Stage 2 將理想權重轉換為可執行交易張數。令 $P_i$ 為每張價格、$x_i^0$ 為期初張數、$C^0$ 為期初現金、$c_b,c_s$ 為買入／賣出交易成本率，且
$
V^0=C^0+\sum_iP_ix_i^0
$
為期初總資產。若股票被選入，最低與最高持有張數分別為
$
L_i^{\text{lot}}=\lceil LV^0/P_i\rceil
$
與
$
U_i^{\text{lot}}=\lfloor UV^0/P_i\rfloor
$；
另定義
$
\bar u_i=\max(0,U_i^{\text{lot}}-x_i^0)
$
為最大可買入張數。

\textbf{Stage 2 模型：}
{\scriptsize
\setlength{\arraycolsep}{3pt}
\begin{align}
\min_{x,u,v,C,B,Q,d}\quad 
& \sum_{i\in I}d_i+\lambda_{\text{trade}}\sum_{i\in I}(B_i+Q_i) \label{obj2}\\[-1mm]
\text{s.t.}\quad
& x_i=x_i^0+u_i-v_i,\quad \forall i\in I, \label{c2:inv}\\
& C=C^0+\sum_{i\in I}P_iv_i(1-c_s)-\sum_{i\in I}P_iu_i(1+c_b), 
&& C\ge0, \label{c2:cash}\\
& d_i\ge \frac{P_ix_i}{V^0}-w_i^\star,\quad \forall i\in I,
&& d_i\ge w_i^\star-\frac{P_ix_i}{V^0},\quad \forall i\in I, \label{c2:dev}\\
& L_i^{\text{lot}}z_i^\star\le x_i\le U_i^{\text{lot}}z_i^\star,\quad \forall i\in I, \label{c2:lot}\\
& u_i\le \bar u_iB_i,\quad v_i\le x_i^0Q_i,\quad \forall i\in I,
&& B_i+Q_i\le1,\quad \forall i\in I, \label{c2:bq}\\
& x_i,u_i,v_i\in\mathbb Z_{\ge0},\quad B_i,Q_i\in\{0,1\},
&& d_i,C\ge0. \nonumber
\end{align}
}

\Cref{c2:inv,c2:cash} 分別表示庫存與現金守恆；\Cref{c2:dev} 線性化實際權重與理想權重的 L1 偏離；\Cref{c2:bq} 限制同一股票不能同時買入與賣出。為維持兩階段一致性，$L_i^{\text{lot}},U_i^{\text{lot}}$ 亦可反向投影回 Stage 1 的權重界。


% ============================================================
\section{Algorithms}
\label{sec-algo}

Stage 1 使用兩個啟發式方法：\textbf{Top-K} 作為下界基準，\textbf{Greedy} 作為主要求解方法。兩者共用四個子程序：\textsc{Allocate}、\textsc{Repair}、\textsc{Diversify} 與 \textsc{LocalSearch}。預處理時先計算每檔股票的有效上限 $\mathit{cap}_i=\min(U_i,\ell_i)$ 與下限 $\mathit{floor}_i=L_i$，並篩出可選股集合。

\textbf{\textsc{Allocate}} 採用 water-filling 配權。給定持股集合 $S$ 後，先將每檔股票配置至 $\mathit{floor}_i$，再依 $\mu_i$ 由高到低分配剩餘權重，同時受限於 $\mathit{cap}_i$、產業剩餘容量與總預算。此步可確保 Budget、Cardinality、Box、Liquidity 與 Industry 限制成立。

\textbf{\textsc{Repair}} 用於修復 Beta 與 CVaR 違反。若 $\beta^\top w>\beta^{\max}$ 或 $\mathrm{CVaR}>\bar C^{\text{risk}}$，演算法會將權重從風險分數較高且仍高於下限的股票，移轉至風險分數較低且仍有容量的股票；若無法透過移轉修復，則進行換股。

\textbf{Top-K} 直接選取 $\mu_i$ 最高的 $K$ 檔股票，先以 \textsc{Allocate} 配權，若因容量或產業限制導致無法投滿，則呼叫 \textsc{Diversify} 進行被迫換股，最後再以 \textsc{Repair} 修復風險限制。Top-K 不進行主動目標改善，因此主要作為比較基準。

\textbf{Greedy} 採用多起點局部搜尋。起點 A 由 industry-aware greedy selection 建構，起點 B 則採用 Top-K 解；兩個起點皆經過 \textsc{Allocate} 與 \textsc{Repair} 後，再執行 1-swap \textsc{LocalSearch}。Local search 會嘗試將既有持股替換為高 $\mu_i$ 的非持股候選，並只接受可行且使 $\mu^\top w$ 嚴格上升的換股。由於 Top-K 是 Greedy 的起點之一，Greedy 的目標值保證不低於 Top-K。

\begin{table}[!hbt]
\centering
\scriptsize
\setlength{\tabcolsep}{4pt}
\renewcommand{\arraystretch}{0.9}
\begin{tabular}{p{0.22\textwidth}p{0.68\textwidth}}
\toprule
Component & Function \\
\midrule
\textsc{Allocate} & Water-filling allocation subject to lower bounds, caps, liquidity and industry limits. \\
\textsc{Repair} & Adjusts weights or swaps names to restore Beta and CVaR feasibility. \\
\textsc{Diversify} & Replaces names when the current set cannot satisfy capacity or industry constraints. \\
\textsc{LocalSearch} & Performs feasible 1-swap improvements and accepts only objective-improving moves. \\
\bottomrule
\end{tabular}
\caption{Summary of heuristic components}
\label{tab:heuristic-components}
\end{table}


% ============================================================
\section{Data collection and generation}
\label{sec-data}

\textbf{真實世界資料。}{}本研究所用之台股歷史價量資料來自 \emph{台灣經濟新報 (TEJ)}的資料庫，經由標準化欄位名稱、剔除指數列、處理停牌補值（forward-fill 上限 3 日）、單位換算，再輸出為 \texttt{ prices.parquet}；候選股池與市值、產業由 \texttt{candidates.parquet}提供；日頻 join 完成的價量／市值／beta\_3m／ROE／營收成長表為 \texttt{params.parquet}；大盤指數收盤獨立存於 \texttt{benchmark.parquet} 作為 $\beta$ 回歸用之 market factor。整個資料管線之間以 parquet 為介面，可被多人平行開發。

\textbf{參數估計 — Module 2 (\texttt{src/params/})。}
\begin{itemizeTight}
    \item \textbf{$\mu_i$（吸引力分數，\texttt{params/mu.py}）}：基準版採 12-1 動能——過去 11 個月對數報酬、排除最近 1 個月以避開短期反轉；可選「動能 + ROE z-score」混合版（\texttt{config.yaml.mu.method=blend}）。資料若可用視窗不足會自動降階到實際可用天數。
    \item \textbf{$\beta_i$（\texttt{params/beta.py}）}：預設以個股對 TAIEX (Y9999) 之過去 120 日對數報酬做滾動 OLS（method=\texttt{ols}），同檔提供 \texttt{tej\_params}（讀 \texttt{params.parquet} 的 \texttt{beta\_3m}）與 \texttt{tej\_raw}（解析原始 \texttt{raw\_beta.csv}）兩個備援。所有來源最後一律 winsorize 到 $[-1, 3]$ 並把缺值補成 $\beta_i = 1.0$。
    \item \textbf{$\ell_i$（流動性上限，\texttt{params/liquidity.py}）}：$\ell_i = \min(U, \rho \cdot ADV_i / V^0)$，其中 $ADV_i$ 為 20 日滾動成交金額平均、$\rho = 0.10$（單檔部位不超過 ADV 之 10\%）。
    \item \textbf{$r_{i,s}$（CVaR 情境，\texttt{params/cvar\_scenarios.py}）}：基準版採歷史情境法——取過去 240 個交易日的日報酬 $(r_t)_{t=t-239}^{t}$ 作為 $|\Omega| = 240$ 個情境，逐日為一個情境向量；報酬以對數型式計算。另提供 moving-block bootstrap（method=\texttt{block\_bootstrap}）作為情境擴增手段，符合開發計畫 §9 對 $|\Omega| \geq 250$ 之要求。
    \item \textbf{$L_i, U_i$（per-stock 權重界）}：根據 $L_i^{\text{lot}}, U_i^{\text{lot}}$ 反向投影回連續權重；當 $U_i^{\text{lot}} < 1$（一張價格已超過 $U V^0$）則該檔自動退出可選集 $z_i = 0$，避免兩階段可行域不一致。
    \item \textbf{$w_i^0$（期初權重）}：首期由設定檔指定（預設全 0、即「全現金起跑」）；rolling 回測時由上一期 Stage 2 輸出之 $x_i, C$ 反推。
\end{itemizeTight}

\textbf{規模掃描子集（real-data subsets）。}{}為觀察演算法在不同 universe 規模下的 gap 與時間表現，我們以「$\mu$ 前 $N$ 名子集合」由真實 instance 切出 $N \in \{50, 200, 500, 1000\}$ 共四個子集（並覆寫 $K$）；每個子集的 $\beta, \ell, r_{i,s}$ 等其他參數與真實 instance 一致。這些屬於「真實資料的縮放」，用於 \cref{sec-eval} 的規模敏感度分析（\cref{tab:bench-sweep}）。

\textbf{Self-generated instances — 27 個結構化合成實例。}{}為符合作業 §5「self-generated instances 為必需」、並系統性檢驗演算法在「真實快照永遠落在可行域內」之外的各種市場結構下之行為，我們另外以一套 10 因子實驗設計（DoE）生成 27 個\emph{完全合成}的 instance。十個因子涵蓋：universe 規模 $|I|\in\{20,50,100\}$、情境數 $|\Omega|\in\{100,250,500\}$、產業形狀、報酬分佈、相關結構、$\beta$ 分佈、$\mu$–$\beta$ 相關、價格分佈（uniform／log-normal，後者刻意觸發 lot-infeasibility）、期初持有、與壓力組合。所有非變因鎖在中心格 baseline，各子集以 OFAT方式單變一因子，便於在報告中逐列解釋；五個子集 S1–S5 共 27 個 instance（\cref{tab:synth-doe}）。價格以單因子 GBM 模擬、$\beta$ 以兩階段「先抽目標 $\beta$ 再餵入模擬器當 loading」構造、$\mu$ 經 Cholesky 與 $\beta$ 耦合以實現 F7；RNG 以 \texttt{numpy.SeedSequence.spawn()} 為各生成模組獨立散播，確保「同 seed $\to$ 四個 parquet 之 bytes 完全相同」（25 個 pytest 驗證）。合成 instance 與真實 pipeline \emph{共用同一套 Stage 1／Stage 2 程式碼}，不更動任何 \texttt{MILP\_phase\_*} 模型檔。三層驗證（schema $\to$ economic $\to$ feasibility）結果為 \textbf{24 OK\_FEASIBLE / 3 OK\_RELAXED}（後者全部落在 S4 壓力子集），27 個 instance 全於 $0.23$\,s 內解出。

\begin{table}[!hbt]
\centering
\small
\begin{tabular}{l r p{0.44\textwidth} p{0.18\textwidth}}
\toprule
子集 & $n$ & 變因（factor $\times$ levels） & 目的 \\
\midrule
S1 規模掃描   & 9 & $|I|\times|\Omega| = \{20,50,100\}\times\{100,250,500\}$ & runtime vs 規模 \\
S2 分佈壓力   & 6 & $\{t,\text{bootstrap}\}\times\{$indep, block, factor$\}$ & 尾部／相關結構 \\
S3 約束壓力   & 5 & wide $\beta$／$\mu$-$\beta=\pm0.5$／集中產業／log-normal 價 & 單一限制 binding \\
S4 不可行陷阱 & 4 & normal／tight CVaR$(0.008)$／tight industry／infeas trap & relaxation 路徑 \\
S5 Warm-start & 3 & S1 對角三格翻 warm（$w^0\neq0$） & turnover binding \\
\bottomrule
\end{tabular}
\caption{自生成合成實例組之 DoE：5 個子集、共 27 個 instance。中心格 baseline 為 $|I|=50,\,|\Omega|=250$、5 balanced 產業、Gaussian、block 相關、tight $\beta$、無 $\mu$-$\beta$ 相關、uniform 價、cold start、normal caps。}
\label{tab:synth-doe}
\end{table}

\textbf{真實 instance 規模。}{}以 2026/05/20 為再平衡日（資料抽取於 2026-05-25 重跑全條 pipeline 後之最新交易日）的真實案例為：1{,}052 檔候選股、34 個產業、240 個歷史情境；其中 46 檔股票因「1 張價格已超過 $U V^0$」（如 2330 台積電、2454 聯發科、3008 大立光）被自動標為 lot-infeasible 並從 $z = 1$ 候選集排除（這是 proposal §10 列出的模型限制 (2)：高股價迫使部分股票無法以單張權重 $\leq U$ 形式持有）。上游候選股池涵蓋 2{,}323 檔（含產業分類，當期 \texttt{factors.parquet} 之 $\mu$ 均值 $0.173$、$\beta$ 均值 $0.440$），在 $\texttt{min\_market\_cap} \geq 5 \times 10^9$ 與 $\texttt{min\_history\_days} \geq 200$ 兩道篩選後得 1{,}052 檔。

\textbf{介面契約（§7）。}{}所有跨模組資料以固定 schema 的 parquet／JSON／CSV 交換：上游 \texttt{prices/candidates/params/benchmark.parquet}（Module 1）、中游 \texttt{factors/scenarios.parquet}（Module 2 於 \texttt{as\_of} 算齊 $(\mu_i,\beta_i,\ell_i,r_{i,s})$）、下游 \texttt{ideal\_portfolio.json}（Stage 1$\to$2）與 \texttt{order\_sheet.csv}（Stage 2$\to$回測／下單）。schema 鎖死於開發起點，使 4 人組得以平行開發。



% ============================================================
\section{Performance evaluation}
\label{sec-eval}

本節在同一份輸入上對照三種求解策略：\textbf{(a) MILP 最佳解}（CBC 後端，MIP gap 容忍度設為 $0$、求解時間上限 $300$\,s）、\textbf{(b) 本研究提出之貪婪啟發式（greedy heuristic）}、與 \textbf{(c) Top-$K$ 等權基準}。最佳化落差（optimality gap）定義為 $\text{gap} = (z^{\text{MILP}} - z) / |z^{\text{MILP}}|$，以百分比表示；目標為極大化，故 gap 為正代表該方法之目標值低於 MILP 最佳解。

\subsection{Stage 1 規模掃描（合成 instances + 真實 instance）}

\Cref{tab:bench-sweep} 為五個規模（$N \in \{50, 200, 500, 1000, 1052\}$）上的對照結果，可歸納三項觀察：
\begin{itemizeTight}
    \item 本研究之貪婪啟發式相對 MILP 之 gap 始終落在 \textbf{0.93\%--1.56\%}，且與 universe 規模無明顯關聯；
    \item Top-$K$ 基準之 gap 隨 universe 規模擴大而\emph{大致單調增加}（$9.80\% \to 15.05\%$）：universe 越大，$\mu$ 排名前 $K$ 名以外的可替換個股越多，單純「依 $\mu$ 排序」所遺漏的機會也越大；此正凸顯啟發式中「換股／局部搜尋（local search）」步驟的價值；
    \item MILP 求解時間自 $N=50$ 的 \textbf{0.21\,s} 近線性成長至 $N=1052$ 的 \textbf{2.19\,s}（CBC 單核），完全落於可實用範圍內。
\end{itemizeTight}

\begin{table}[!hbt]
\centering
\small
\begin{tabular}{rrl rrrrr}
\toprule
$N$ & $K$ & method & objective & gap (\%) & $\beta$ & CVaR & time (s) \\
\midrule
\multirow{3}{*}{50} & \multirow{3}{*}{15}
   & MILP optimal      & 2.3074 & 0.00 & 0.874 & 0.0400 & 0.21 \\
&  & greedy heuristic  & 2.2715 & 1.56 & 0.754 & 0.0400 & 3.59 \\
&  & Top-$K$ baseline  & 2.0813 & 9.80 & 0.790 & 0.0399 & 0.01 \\
\midrule
\multirow{3}{*}{200} & \multirow{3}{*}{20}
   & MILP optimal      & 2.3085 & 0.00 & 0.799 & 0.0400 & 0.50 \\
&  & greedy heuristic  & 2.2869 & 0.93 & 0.837 & 0.0400 & 4.69 \\
&  & Top-$K$ baseline  & 1.9966 & 13.51 & 0.709 & 0.0400 & 0.01 \\
\midrule
\multirow{3}{*}{500} & \multirow{3}{*}{20}
   & MILP optimal      & 2.3085 & 0.00 & 0.799 & 0.0400 & 1.11 \\
&  & greedy heuristic  & 2.2847 & 1.03 & 0.803 & 0.0400 & 6.01 \\
&  & Top-$K$ baseline  & 1.9999 & 13.37 & 0.772 & 0.0400 & 0.01 \\
\midrule
\multirow{3}{*}{1000} & \multirow{3}{*}{20}
   & MILP optimal      & 2.3085 & 0.00 & 0.799 & 0.0400 & 2.43 \\
&  & greedy heuristic  & 2.2842 & 1.05 & 0.793 & 0.0400 & 4.93 \\
&  & Top-$K$ baseline  & 1.9611 & 15.05 & 0.763 & 0.0399 & 0.01 \\
\midrule
\multirow{3}{*}{\textbf{1052}} & \multirow{3}{*}{\textbf{20}}
   & MILP optimal      & \textbf{2.3085} & \textbf{0.00} & 0.799 & 0.0400 & \textbf{2.19} \\
&  & greedy heuristic  & 2.2842 & 1.05 & 0.793 & 0.0400 & 4.92 \\
&  & Top-$K$ baseline  & 1.9611 & 15.05 & 0.763 & 0.0399 & 0.01 \\
\bottomrule
\end{tabular}
\caption{Stage 1 三種方法在不同規模 universe 上的 objective、gap、與時間（rebalance date = 2026/05/20、$|\Omega|=240$、CBC backend）。最後一列為真實 instance；其餘為 self-generated 子集合。}
\label{tab:bench-sweep}
\end{table}

\subsection{Stage 1 合成實例 DoE 對照（27 個 self-generated instance）}

\Cref{tab:synth-bench} 為三種方法（MILP optimal、greedy heuristic、Top-$K$ baseline）在 27 個合成 instance 上、依五個子集平均之最佳化落差，可得三項結論：
\begin{itemizeTight}
    \item 在常規可行的子集（S1--S3）上，貪婪啟發式幾乎完全貼合 MILP 最佳解（平均 gap $\approx 0.00\%$），Top-$K$ 則落後 $0.3\%$--$1.6\%$；全 27 個 instance 中，啟發式\textbf{於 26 個上擊敗 Top-$K$}，唯一例外為 small-universe 結合 warm-start 的邊角案例 SYN\_025。
    \item S4「不可行陷阱」子集中，啟發式與 Top-$K$ 均出現約 $32\%$ 的大幅落差——因 MILP 係透過放寬機制（relaxation fallback）跳離原始模型方求得解，而啟發式受限於\emph{原始}約束、不含放寬機制；此對比恰好凸顯兩者的設計權衡。
    \item S5「warm-start」子集中，啟發式平均略\emph{勝} MILP（$-2.50\%$）：在部分 seed 下周轉率限制未綁定，啟發式得以於較寬鬆的解空間中找到更高目標值，而 MILP 則在 MIP gap 容忍度下提前停止。
\end{itemizeTight}

\begin{table}[!hbt]
\centering
\small
\begin{tabular}{l r r r r}
\toprule
子集 & $n$ & mean opt.\ obj & heuristic gap & Top-$K$ gap \\
\midrule
S1 規模掃描   & 9 & 1.18 & 0.00\% & 0.42\% \\
S2 分佈壓力   & 6 & 1.32 & 0.00\% & 0.28\% \\
S3 約束壓力   & 5 & 1.21 & 0.00\% & 1.61\% \\
S4 不可行陷阱 & 4 & 1.30 & 31.87\% & 31.97\% \\
S5 Warm-start & 3 & 1.14 & $-2.50\%$ & $-4.30\%$ \\
\bottomrule
\end{tabular}
\caption{三種方法在 27 個合成 instance 上之平均最佳化落差（相對 MILP optimal）。S4 之大 gap 源於 MILP 經放寬機制跳離原模型；S5 之負 gap 源於部分 seed 下周轉率未綁定。啟發式整體以 26/27 擊敗 Top-$K$。}
\label{tab:synth-bench}
\end{table}

\textbf{放寬機制（relaxation fallback）之實證。}{}真實 instance 全程未觸發任何放寬，其唯二作用中（active）限制為 CVaR 與周轉率；合成實例的 S4 子集則正面驗證了 \cref{sec-model} 所設計的逐步放寬程序：SYN\_022（CVaR 上限緊縮至 $0.008$）僅放寬 CVaR 一層即恢復可行，SYN\_023（產業上限緊縮）與 SYN\_024（不可行陷阱）則完整走過 CVaR $\to$ 周轉率 $\to$ Beta $\to$ 產業 四層放寬。此結果證明 \cref{sec-model} 所訂之放寬順序確實按設計運作，且僅在原模型不可行時才啟動。

\subsection{Stage 2 整數張數落地 — 真實案例}

我們以部署的 Stage 1 最佳解（目標值 $2.3079$、$\beta = 0.84$，與 \cref{tab:bench-sweep} 之嚴格最佳上界 $2.3085$ 相差 $<0.03\%$）為輸入，在真實 instance 上求解 Stage 2，參數設定為初始資金 $V^0 = 10{,}000{,}000$ 元、買進成本率 $c_b = 0.001425$、賣出成本率 $c_s = 0.004425$、交易懲罰係數 $\lambda_{\text{trade}} = 0.0005$。求解狀態為 Optimal；總 L1 偏離 $\sum_i d_i = \textbf{0.0366}$（即實際相對權重相對理想平均偏離 $0.18\%$）；交易筆數為 $20$，全為買進、無賣出（因首期 $x^0 \equiv 0$）；剩餘現金 \textbf{23{,}804.15} 元，佔 $V^0$ 之 $0.238\%$。20 檔持股中有 13 檔達 $d_i \approx 0$，僅 7 檔因張數離散粒度而偏離。\Cref{tab:order-sheet} 列出偏離最大的五檔，用以說明 Stage 2 何以並非單純的連續權重取整：其偏離全數源於高股價個股無法在單檔上限 $U V^0 = 1{,}000{,}000$ 元的預算內補足權重，或低股價個股「每張權重 $= P_i / V^0$」之離散粒度過粗。

\begin{table}[!hbt]
\centering
\small
\begin{tabular}{l r r r r r p{4.5cm}}
\toprule
stock\_id & 每張價格 & 買數 & w\_actual & w\_ideal & deviation & 備註 \\
\midrule
6727 亞泰金屬 & 496{,}500 & 2 & 0.0993 & 0.0829 & +0.0164 & 一張 $\approx 4.97\% V^0$，2 張即接近 $U$ 上限 \\
6990 華鉬     & 138{,}000 & 4 & 0.0552 & 0.0623 & $-$0.0071 & 5 張會超過流動性上限 $\ell_i$ \\
3054 立萬利   &  66{,}600 & 3 & 0.0200 & 0.0242 & $-$0.0042 & 每張 0.67\%，4 張會超過產業上限 \\
4577 達航科技 & 118{,}000 & 1 & 0.0118 & 0.0153 & $-$0.0035 & 一張 $\approx 1.18\% V^0$ \\
3555 博士旺   & 195{,}000 & 4 & 0.0780 & 0.0807 & $-$0.0027 & 一張 $\approx 1.95\% V^0$ \\
其他 15 檔    & — & — & — & — & $\leq 0.003$ & 含 13 檔 $d_i \approx 0$ \\
\bottomrule
\end{tabular}
\caption{Stage 2 偏離最大的 5 檔（按 $|d_i|$ 排序）；剩餘現金 23{,}804.15 元（0.238\%）。}
\label{tab:order-sheet}
\end{table}

\subsection{研究問題回答}

我們以前述實證結果，逐一回應本研究於計畫階段所提出的五個研究問題：
\begin{itemize}[leftmargin=2em]
    \item \textbf{RQ1：受限最佳化是否優於 Top-$K$ 等權？} 是。在真實 instance 上 Stage 1 MILP 之 $\mu^\top w$ 比 Top-$K$ 高 \textbf{17.71\%}（2.3085 vs.\ 1.9611）；在 5 個規模上皆嚴格 dominate。
    \item \textbf{RQ2：風險、產業與周轉限制是否提升穩定性？} 在部署的 Stage 1 解中 CVaR $= \bar{C}^{\text{risk}} = 0.04$ 與 turnover $= \tau = 1.00$（首期 $w^0 \equiv 0$ 之結構性下限）兩條限制 active；組合 $\beta = 0.84$ 未觸頂於 $\beta^{\max} = 1.10$，顯示在當期動能驅動下，下檔風險與買進預算才是主導限制。
    \item \textbf{RQ3：連續權重轉整數張數的偏離有多大？} 累積 $\sum_i d_i = 0.0366$（3.66\%）；20 檔中 13 檔完全對齊，僅 7 檔因離散粒度而偏離；最大單檔偏離 0.0164 來自一張價格約 50 萬的 6727 亞泰金屬。
    \item \textbf{RQ4：納入現金守恆與交易成本後是否仍能執行？} 是。剩餘現金 23{,}804 元（0.238\%）、$C \geq 0$ 始終成立。
    \item \textbf{RQ5：兩階段相較單階段是否更貼近真實流程？} 是。Stage 2 額外處理了單階段連續權重模型完全忽略的「整數張數、現金守恆、買賣互斥、交易手續費」四項實務限制，所產生的下單清單可直接交付執行；反之，單階段模型輸出的連續權重仍須人工事後取整，容易引入二次偏離。
\end{itemize}



% ============================================================
\section{Conclusions}
\label{sec-conclusion}

本研究將「主題型台股投資的再平衡決策」嚴謹地形式化為兩個串接的 MILP 子問題，並以真實 TEJ 資料與自我生成的規模掃描 instance 雙軌驗證整體系統。Stage 1 求出同時考量 Beta、CVaR、產業集中度、周轉率與流動性的理想連續權重；Stage 2 再將其無縫轉化為符合整數張數、現金守恆與買賣互斥的可執行下單清單。此外，我們提出一套多起點貪婪啟發式作為自有演算法：在 1{,}052 檔 universe 上其最佳化落差僅 \textbf{1.05\%}，顯著優於 Top-$K$ 等權基準（gap $15.05\%$），且無需商用求解器即可於秒級內給出可行解。值得注意的是，即使規模上探至 $N = 1052$，MILP 仍能於約 $2.2$ 秒內求得最佳解；因此該啟發式於現階段並非必要，但作為求解失敗時的後備方案、以及未來向更大規模問題（如多期聯立最佳化）延伸的基礎工具，仍具實用價值。

\textbf{後續工作。}{}本研究尚有以下可延伸方向：
\begin{itemizeTight}
    \item 滾動式回測（rolling backtest）：以本兩階段系統、Top-$K$ 等權、與純 Stage 1 連續權重三條策略，進行 5 年月頻／季頻回測，比較 NAV、Sharpe 比率、最大回撤（max drawdown）、實際周轉率與現金占比。
    \item 將吸引力分數延伸為多因子或機器學習預測（以 walk-forward 防止前視偏誤）；
    \item 引入更細緻的交易衝擊成本函數（如 sq-root law）；
    \item 將整體模型延伸為真正的 multi-period optimization（目前每期獨立求解，未捕捉跨期再平衡的耦合）。
\end{itemizeTight}

\end{document}
